\chapter{مبانی نظری معادلات غیرخطی و تقریب خطی}

\section{مقدمه}
بسیاری از مسئله‌های ریاضی، فیزیکی و مهندسی در نهایت به حل یک معادله غیرخطی می‌رسند. برای نمونه، اگر رابطه بین متغیرها به صورت $f(x)=0$ نوشته شود، هدف یافتن مقداری از $x$ است که تابع را صفر کند. در حالت کلی، چنین مقداری به صورت بسته و دقیق به دست نمی‌آید؛ بنابراین به جای ریشه دقیق، دنباله‌ای از تقریب‌ها ساخته می‌شود. روش نیوتن یکی از مهم‌ترین روش‌های تکراری برای این کار است و در تحلیل عددی جایگاه ویژه‌ای دارد \cite{burden,atkinson}.

در این گزارش، ایده اصلی روش نیوتن، شرط‌های همگرایی و یک مثال عددی بررسی می‌شود. دلیل انتخاب این موضوع آن است که هم از نظر ریاضی پر از فرمول‌های قابل ارجاع است و هم از نظر الگوریتمی، پیاده‌سازی ساده و قابل فهمی دارد.

\section{تعریف ریشه و خطای تقریب}
\begin{definition}[ریشه تابع]
فرض کنید $f:I\subseteq \mathbb{R}\to \mathbb{R}$ تابعی حقیقی باشد. عدد $\alpha\in I$ را ریشه تابع $f$ می‌نامیم، هرگاه
\begin{equation}
 f(\alpha)=0.
 \label{eq:root-definition}
\end{equation}
اگر $f'(\alpha)\neq 0$ باشد، ریشه را ساده می‌نامیم.
\end{definition}

در روش‌های عددی، به جای $\alpha$ معمولاً دنباله $\{x_k\}_{k=0}^{\infty}$ تولید می‌شود. خطای گام $k$ام به صورت زیر تعریف می‌شود:
\begin{equation}
 e_k=x_k-\alpha, \qquad |e_k|=|x_k-\alpha|.
 \label{eq:error-definition}
\end{equation}
هرچه $|e_k|$ کوچک‌تر شود، تقریب $x_k$ به ریشه دقیق نزدیک‌تر است. در عمل چون $\alpha$ معمولاً معلوم نیست، از معیارهای توقف زیر استفاده می‌کنیم:
\begin{equation}
 |x_{k+1}-x_k|<\varepsilon
 \quad \text{یا} \quad
 |f(x_{k+1})|<\varepsilon.
 \label{eq:stopping}
\end{equation}

\section{بسط تیلور و تقریب خطی}
ایده اصلی روش نیوتن بر پایه تقریب خطی تابع است. اگر تابع $f$ در همسایگی نقطه $x_k$ دوبار مشتق‌پذیر باشد، بسط تیلور حول $x_k$ برابر است با
\begin{equation}
 f(x_k+h)=f(x_k)+h f'(x_k)+\frac{h^2}{2}f''(\xi_k),
 \label{eq:taylor}
\end{equation}
که در آن $\xi_k$ بین $x_k$ و $x_k+h$ قرار دارد. اگر جمله درجه دوم را در \cref{eq:taylor} نادیده بگیریم، تقریب خطی زیر به دست می‌آید:
\begin{equation}
 f(x_k+h)\approx f(x_k)+h f'(x_k).
 \label{eq:linearization}
\end{equation}
اکنون اگر بخواهیم نقطه بعدی روی همین تقریب خطی ریشه باشد، کافی است در \cref{eq:linearization} مقدار $f(x_k+h)$ را صفر بگیریم. در نتیجه
\begin{equation}
 0=f(x_k)+h f'(x_k)
 \quad \Longrightarrow \quad
 h=-\frac{f(x_k)}{f'(x_k)}.
 \label{eq:h-formula}
\end{equation}
پس نقطه بعدی برابر با $x_{k+1}=x_k+h$ خواهد بود. این رابطه در فصل بعد به عنوان فرمول اصلی روش نیوتن استفاده می‌شود.

\section{مرتبه همگرایی}
\begin{definition}[مرتبه همگرایی]
فرض کنید $x_k\to \alpha$ و $e_k=x_k-\alpha$. اگر عددهای $p>0$ و $C>0$ وجود داشته باشند به طوری که
\begin{equation}
 \lim_{k\to\infty}\frac{|e_{k+1}|}{|e_k|^p}=C,
 \label{eq:order-conv}
\end{equation}
آنگاه می‌گوییم دنباله $\{x_k\}$ با مرتبه $p$ به $\alpha$ همگرا می‌شود. در حالت $p=2$، همگرایی درجه دوم نامیده می‌شود.
\end{definition}

اهمیت روش نیوتن در این است که تحت فرض‌های مناسب، همگرایی آن درجه دوم است؛ یعنی تقریباً تعداد رقم‌های درست در هر گام دو برابر می‌شود. در \cref{thm:quadratic} این گزاره به صورت دقیق‌تر بیان خواهد شد.
